\SourcePage{282}{285}
\section{Exchange interaction}

The omission of particle spin from the Schrödinger equation in no way diminishes either that equation or any results obtained from it. The reason is that the electric interaction between particles is independent of their spins\footnote{This holds only within the nonrelativistic approximation. Once relativistic effects are included, the interaction of charged particles does depend on spin.}.

In mathematical terms, the Hamiltonian of a system of electrically interacting particles (with no magnetic field present) contains no spin operators and therefore acts on a wave function without affecting its spin variables. In fact, each component of the wave function consequently satisfies the Schrödinger equation. Equivalently, the wave function of a particle system may be written as the product
\[
 \psi(\xi_1,\xi_2)=\chi(\sigma_1,\sigma_2,\ldots)\varphi(\mathbf r_1,\mathbf r_2,\ldots),
\]
where $\varphi$ depends only on the particle coordinates, and $\chi$ only on their spins. We shall refer to the former as the \emph{coordinate} or \emph{orbital} wave function and to the latter as the \emph{spin} wave function.

The Schrödinger equation essentially determines only the coordinate function $\varphi$, leaving $\chi$ arbitrary. Whenever the particle spins themselves are of no interest, one may therefore apply the Schrödinger equation with the coordinate function alone as the wave function, as was done in the preceding chapters.

Nevertheless, although the electric interaction is independent of spin in the manner just described, the system's energy has a peculiar dependence on its total spin, arising ultimately from the indistinguishability principle for identical particles.

Consider a system of just two identical particles. Solving the Schrödinger equation yields a sequence of energy levels, each associated with a definite symmetric or antisymmetric coordinate wave function $\varphi(\mathbf r_1,\mathbf r_2)$.

Indeed, particle identity makes the Hamiltonian of the system, and therefore its Schrödinger equation, invariant under interchange of the particles. If an energy level is nondegenerate, exchanging the coordinates $\mathbf r_1$ and $\mathbf r_2$ can change $\varphi(\mathbf r_1,\mathbf r_2)$ only

\SourcePage{283}{286}
by a constant factor; a second exchange shows that this factor can only be $\pm1$\footnote{In a degenerate level, linear combinations of its functions can always be chosen so as to satisfy this condition as well.}.

Suppose first that the particles have zero spin. There is then no spin factor at all, and the wave function reduces to the coordinate function $\varphi(\mathbf r_1,\mathbf r_2)$ alone. This function must be symmetric, since zero-spin particles obey Bose statistics. Hence not every energy level produced by a formal solution of the Schrödinger equation can actually occur: levels associated with antisymmetric functions $\varphi$ are impossible for the system under consideration.

Interchanging two identical particles is equivalent to inverting the coordinate system, whose origin is chosen at the midpoint of the line joining them. Under inversion, on the other hand, $\varphi$ must acquire the factor $(-1)^l$, where $l$ is the orbital angular momentum of the relative motion of the two particles (see §~30). Combining this fact with the preceding result, we conclude that a pair of identical zero-spin particles can have only even orbital angular momentum.

Next let the system consist of two spin-$1/2$ particles, electrons for example. The full wave function---the product of $\varphi(\mathbf r_1,\mathbf r_2)$ and the spin function $\chi(\sigma_1,\sigma_2)$---must necessarily be antisymmetric under interchange of the two particles. A symmetric coordinate function must therefore be accompanied by an antisymmetric spin function, and conversely. Write the spin function in spinor form as a second-rank spinor $\chi^{\lambda\mu}$, with each index corresponding to the spin of one electron. A function symmetric in the spins is represented by a symmetric spinor ($\chi^{\lambda\mu}=\chi^{\mu\lambda}$), and an antisymmetric function by an antisymmetric spinor ($\chi^{\lambda\mu}=-\chi^{\mu\lambda}$). But a symmetric second-rank spinor, as we know, describes a system of total spin one, whereas an antisymmetric spinor reduces to a scalar and thus corresponds to spin zero.

We arrive at the following result. Energy levels whose Schrödinger-equation solutions $\varphi(\mathbf r_1,\mathbf r_2)$ are symmetric can actually occur

\SourcePage{284}{287}
when the total spin of the system is zero, that is, when the two electron spins are ``antiparallel'' and add to zero. Energy values associated with antisymmetric functions $\varphi(\mathbf r_1,\mathbf r_2)$ require total spin one, so the electron spins must be ``parallel.''

In other words, the possible energies of an electron system depend on its total spin. One may on this basis speak of a special interaction between the particles that produces this dependence. It is called the \emph{exchange interaction}. This is a purely quantum effect and disappears completely, along with spin itself, in the limiting passage to classical mechanics.

The two-electron system just examined has this characteristic feature: every energy level corresponds to one definite total-spin value, either 0 or 1. As we shall see below (§~63), this one-to-one correspondence between energy levels and spin values persists in systems containing any number of electrons. It does not, however, hold in systems made up of particles whose spin exceeds $1/2$.

Consider two particles of arbitrary spin $s$. Their spin wave function is a spinor of rank $4s$,
\[
 \chi^{\overbrace{\lambda\mu\ldots}^{2s}\overbrace{\rho\sigma\ldots}^{2s}},
\]
one half ($2s$) of whose indices correspond to the spin of one particle, and the other half to that of the other. The spinor is symmetric in the indices within each group. Interchanging the particles exchanges every index $\lambda,\mu,\ldots$ in the first group with the indices $\rho,\sigma,\ldots$ in the second. To obtain the spin function for a system state of total spin $S$, this spinor must be contracted in $2s-S$ pairs of indices, each pair comprising one index from each group, and symmetrized over the remaining indices. The result is a symmetric spinor of rank $2S$.

As we know, contracting a spinor in a pair of indices means forming a combination antisymmetric in those indices. Thus interchange of the particles multiplies the spin wave function by $(-1)^{2s-S}$.

The full wave function of the two-particle system, on the other hand, must acquire the factor $(-1)^{2s}$ upon particle interchange: $+1$ for integral $s$, and $-1$ for half-integral $s$. It follows that the symmetry of the coordinate wave function under particle interchange is determined by the factor $(-1)^S$, which depends only on $S$.

\SourcePage{285}{288}
We have therefore shown that the coordinate wave function for two identical particles is symmetric when their total spin is even and antisymmetric when it is odd.

Recalling the relation established above between particle interchange and inversion of the coordinate system, we also conclude that for even (odd) $S$ the system can possess only even (odd) orbital angular momentum.

Here too, then, the system's possible energies display a dependence on total spin, but the relation is no longer unique. Energy levels associated with symmetric (antisymmetric) coordinate wave functions may occur for every even (odd) value of $S$.

Let us count all distinct states of a two-particle system having even and odd values of $S$. The quantity $S$ takes $2s+1$ values: $2s,2s-1,\ldots,0$. For each fixed $S$ there are $2S+1$ states, distinguished by the $z$ component of spin, giving $(2s+1)^2$ states altogether. Let $s$ be integral. Of the $2s+1$ possible values of $S$, $s+1$ are even and $s$ are odd. The total number of states with even $S$ is
\[
 \sum_{S=0,2,\ldots,2s}(2S+1)=(2s+1)(s+1);
\]
the remaining $s(2s+1)$ states have odd $S$. Similarly, for half-integral $s$ there are $s(2s+1)$ states with even $S$ and $(s+1)(2s+1)$ with odd $S$.

\ProblemHeading

\textbf{1.} Determine the exchange splitting of the energy levels of a two-electron system, treating the interaction between the electrons as a perturbation.

\SolutionHeading Let the particles, with their interaction omitted, occupy states with orbital wave functions $\varphi_1(\mathbf r_1)$ and $\varphi_2(\mathbf r_2)$. The states of total spin $S=0$ and $S=1$ correspond, respectively, to the symmetrized and antisymmetrized products
\[
 \varphi=\frac1{\sqrt2}[\varphi_1(\mathbf r_1)\varphi_2(\mathbf r_2)\pm\varphi_1(\mathbf r_2)\varphi_2(\mathbf r_1)].
\]
In these states the mean values of the particle-interaction operator $U(\mathbf r_2-\mathbf r_1)$ are $A\pm J$, where
\[
 A=\iint U|\varphi_1(\mathbf r_1)|^2|\varphi_2(\mathbf r_2)|^2,dV_1dV_2,
\]
\[
 J=\iint U\varphi_1(\mathbf r_1)\varphi_1^*(\mathbf r_2)\varphi_2(\mathbf r_2)\varphi_2^*(\mathbf r_1),dV_1dV_2
\]
(the integral $J$ is called the \emph{exchange integral}). Dropping the additive constant $A$, which is not of exchange character, gives the level shifts $\Delta E_0=J$ and $\Delta E_1=-J$; the subscript specifies $S$. These quantities

\SourcePage{286}{289}
may be represented as the eigenvalues of the spin \emph{exchange operator}\footnote{This operator was introduced by Dirac.}
\begin{equation}
 \widehat V_{\text{ex}}=-\frac12J(1+4\hat{\mathbf s}_1\hat{\mathbf s}_2)
 \tag{1}
\end{equation}
(for the eigenvalues of the product $\mathbf s_1\mathbf s_2$, see problem 2 of §~55).

If, for example, the two electrons belong to different atoms, the exchange integral decreases exponentially as the interatomic distance $R$ increases. The form of the integrand makes clear that this integral is governed by the ``overlap'' of the state wave functions $\varphi_1(\mathbf r_1)$ and $\varphi_2(\mathbf r_2)$. Using the asymptotic decay law for discrete-spectrum wave functions (cf. (21.6)), one finds
\[
 J\sim\exp[-(\kappa_1+\kappa_2)R],\qquad
 \kappa_1=\frac1\hbar\sqrt{2m|E_1|},\qquad
 \kappa_2=\frac1\hbar\sqrt{2m|E_2|},
\]
where $E_1,E_2$ are the electron energy levels in the two atoms.

\textbf{2.} Do the same for a system of three electrons.

\SolutionHeading From formula (1) of problem 1, the operator for the pairwise exchange interaction of three electrons is
\begin{equation}
 \widehat V_{\text{ex}}=-\sum J_{ab}(1/2+2\hat{\mathbf s}_a\hat{\mathbf s}_b),
 \tag{1}
\end{equation}
where the sum is over the particle pairs 12, 13, and 23. The matrix elements of $\hat{\mathbf s}_a\hat{\mathbf s}_b$ between states with different pairs of values $\sigma_a,\sigma_b$ follow from (55.6) and are
\[
 \langle1/2,1/2|\mathbf s_a\mathbf s_b|1/2,1/2\rangle=1/4,
 \quad \langle1/2,-1/2|\mathbf s_a\mathbf s_b|1/2,-1/2\rangle=-1/4,
\]
\[
 \langle1/2,-1/2|\mathbf s_a\mathbf s_b|-1/2,1/2\rangle=1/2.
\]
First determine the energy corresponding to the largest possible projection of the total spin, $M_S=\sigma_1+\sigma_2+\sigma_3$, namely $M_S=3/2$. This also determines the energy of the state with total spin $S=3/2$. Evaluation of the corresponding diagonal matrix element of (1) gives
\[
 \Delta E_{3/2}=-(J_{12}+J_{13}+J_{23}).
\]

Now turn to the states with $M_S=1/2$. This value can be realized in three ways, according to which of $\sigma_1,\sigma_2,\sigma_3$ equals $-1/2$, the other two being $1/2$. These states would therefore lead to a cubic secular equation. The calculation can be shortened at once by observing that one root must correspond to the energy already found for the $S=3/2$ state. The secular equation must consequently be divisible by $\Delta E-\Delta E_{3/2}$, allowing us in this instance to avoid calculating the constant term of the cubic\footnote{This device is particularly useful in analogous calculations for systems with more particles.}.

Specifically, calculation of the leading terms gives
\[
 (\Delta E)^3+(J_{12}+J_{13}+J_{23})(\Delta E)^2+
 [J_{12}J_{13}+J_{12}J_{23}+J_{13}J_{23}-(J_{12}^2+J_{13}^2+J_{23}^2)]\Delta E+\ldots=0,
\]
and division by $\Delta E+J_{12}+J_{13}+J_{23}$ yields the two energy levels belonging to

\SourcePage{287}{290}
states of spin $S=1/2$:
\[
 \Delta E_{1/2}=\pm[J_{12}^2+J_{13}^2+J_{23}^2-J_{12}J_{13}-J_{12}J_{23}-J_{13}J_{23}]^{1/2}.
\]
There are thus three energy levels in all, in agreement with the count made in problem 1 of §~63.

\textbf{3.} From which states can the ${}^{8}\mathrm{Be}$ nucleus decay into two $\alpha$ particles?

\SolutionHeading Since an $\alpha$ particle has no spin, a system of two $\alpha$ particles can have only even orbital angular momentum (equal to its total angular momentum), and its states have even parity. The stated decay can therefore occur only from even-parity states of ${}^{8}\mathrm{Be}$ having even total angular momentum.
